\documentclass{article}

\def\CourseCode{ECEN-662}
\def\CourseName{Estimation \& Detection Theory}
\def\Instructor{Dr. Tie Liu}
<<<<<<< HEAD
\def\Semester{Spring 2026}
\def\Student{Brody Jordan}
\def\University{Texas A\&M University}
\def\Cover{figures/_cover.pdf}
=======
\def\Semester{Spring 2026}
>>>>>>> 249e4c317b2a4234368ca0c3c75e714883e8075a

\usepackage[
    AssignmentName={Homework 7},
    DueDate={April 12, 2026 at 11:59 PM}
]{assignment}


\usepackage{cancel}

\begin{document}
\MakeAssignmentTitle

% ------ PROBLEM 1 ------ %
\begin{homeworkProblem}

  Given a uniform random variable $U \sim \Uc([0,1])$, derive the inverse transform sampling function required
  to generate samples from a Rayleigh distribution with parameter $\sigma$:
  \[
    f(x) = \frac{x}{\sigma^2} \exp \bp{- \frac{x^2}{2\sigma^2}}, \quad x \geq 0.
  \]

  \solution


\end{homeworkProblem}

\pagebreak

% ------ PROBLEM 2 ------ %
\begin{homeworkProblem}

  Suppose we want to sample from a standard Normal distribution, but only the positive half (a truncated normal).
  The target density
  \[
    f(x) \propto \exp (-x^2/2), \quad x > 0.
  \]
  We will use an exponential distribution
  \[
    g(x) = \lambda \exp(-\lambda x), \quad x \geq 0
  \]
  as our proposal distribution.
  \begin{enumerate}[label=(\roman*), topsep=5pt, itemsep=5pt]
    \item Find the optimal scaling constant $M$ as a function of $\lambda$ such that $f(x) \leq M g(x)$ for all $x \geq 0$.
    \item Analytically determine the optimal rate parameter $\lambda$ that minimizes the rejection rate. \newline
  \end{enumerate}

  \solution

\end{homeworkProblem}

\pagebreak

% ------ PROBLEM 3 ------ %
\begin{homeworkProblem}
  We want to estimate the probability of a standard normal random variable exceeding a threshold of 4:
  \[
    \P(X > 4) = \int_4^\infty \frac{1}{\sqrt{2\pi}} \exp\bp{-\frac{x^2}{2}} \, dx.
  \]
  \begin{enumerate}[label=(\roman*), topsep=5pt, itemsep=5pt]
    \item If we use standard Monte Carlo by drawing $N = 10^5$ samples from $X \sim \Nc(0, 1)$, what is the expected
          number of samples that will fall in the region $X > 4$?
    \item Propose a shifted Gaussian importance sampling distribution
          \[
            q(x) = \Nc(\mu, \sigma^2)
          \]
          that would significantly reduce the variance of the estimate. Justify your choice of $\mu$ and write out
          the final Importance Sampling estimator equation. \newline
  \end{enumerate}

  \solution

\end{homeworkProblem}

\pagebreak

% ------ PROBLEM 4 ------ %
\begin{homeworkProblem}

  Let the target distribution be an unnormalized symmetric bimodal distribution:
  \[
    \pi(x) \propto \exp \bp{-\frac{(x-3)^2}{2}} + \exp\bp{-\frac{(x+3)^2}{2}}.
  \]
  We will use a Metropolis-Hastings algorithm with a Gaussian random walk proposal:
  \[
    p(x, x^{\prime}) = \Nc(x^{\prime}; x, \sigma_p^2).
  \]
  Assume the current state of the chain is $x_t = 3.0$ (sitting precisely on the right peak).
  \begin{enumerate}[label=(\roman*), topsep=5pt, itemsep=5pt]
    \item Calculate the exact acceptance ratio $\alpha(x_t, x^{\prime})$ if the proposal draws a new state $x^{\prime} = 0.0$ (the deep valley between the peaks).
    \item Conceptually, what happens to the Markov Chain if the proposal variance $\sigma_p^2$ is set extremely
          small (e.g., 0.01)? What happens if it is set extremely large (e.g., 100)? \newline
  \end{enumerate}

  \solution

\end{homeworkProblem}

\end{document}

